\documentclass[10pt,a4paper]{amsart}
\usepackage[margin=19mm]{geometry}
\usepackage{amsmath,amssymb,mathtools}
\usepackage[scr=rsfs,cal=euler]{mathalfa}
\usepackage{xfrac}
\usepackage[unicode,hidelinks,bookmarks=false]{hyperref}
\newcommand{\ie}{i.e.\ }
\newcommand{\cf}{cf.\ }
\newcommand{\resp}{respectively\ }
\newcommand{\Subsection}{\subsection}
\providecommand{\nobreakdash}{\nobreak\hbox{-}}
\setlength{\emergencystretch}{2em}
\multlinegap0pt
\usepackage{amssymb}
\usepackage{float}
\usepackage{braket}
\newtheorem{theorem}{Theorem}
\title{Optimal graphons in the edge-2star model}
\author{Charles Radin, Lorenzo Sadun}
\date{Source: arXiv:2305.00333v1 (CC BY 4.0)}
\begin{document}
\maketitle

\noindent\emph{Source and licence.} Optimal graphons in the edge-2star model. Version arXiv:2305.00333v1 (CC BY 4.0), distributed by arXiv under the Creative Commons Attribution 4.0 International licence, http://creativecommons.org/licenses/by/4.0/, as recorded in the arXiv licence metadata for this exact version and shown on the abstract page https://arxiv.org/abs/2305.00333v1. Full licence notice: \texttt{licenses/arxiv-2305.00333-CC-BY-4.0.txt}.\par\smallskip

\noindent\textit{Editorial scope.} Counted unit: the article's theorem thm:low-t (source0.tex lines 683--695). The article's complete original proof of it is delivered with it (source0.tex lines 714--933, one \texttt{proof} environment of 220 lines, which the article places away from the statement). No mathematics is written, repaired or re-derived: the statement and the proof are the article's own lines, in the article's own notation, and no retained line is substituted or re-flowed.  Retained uncounted as the source-local context the proof needs: lines 465--467; lines 474--476; lines 664--675.  Nothing was omitted from any retained span, so the package carries no omission record.  No retained line contains a citation, so the wrapper carries no bibliography and no bibliography entry.  No retained line contains a figure environment, an image-inclusion command or drawing code, so no image is delivered and the package declares no asset dependency.  Editorial formatting only, in addition: an AMS \texttt{amsart} preamble that restates the article's own declarations and macros used by the retained lines, namely the environments \texttt{theorem} on separate counters, together with the standard packages those lines need.  The retained environments print in this document as Theorem 1 in order of appearance, whereas the article prints the same items with the numbering of its own full document.  The complete native source of the article as distributed in the arXiv e-print for this exact version is bundled unchanged as \texttt{sources/140766/source0.tex}.
\begin{equation}\label{eq1} 
g(x,y) = \frac{1}{1 + \exp(-(\alpha+\beta (d(x)+d(y))))},
\end{equation} 

\begin{equation}\label{eq3} 
z = k(z) : = \int_0^1 \frac{dy} {1 + \exp(-(\alpha+\beta (z+d(y))))}. 
\end{equation}

Let 
\begin{equation} 
\zeta = \sqrt{\tilde t + (e - \frac12)^2}, 
\end{equation}
and consider the bipodal graphon $g_0$ with 
\begin{eqnarray}\label{eq:ansatz}
a & = & 1 - e - 2\zeta, \cr 
b & = & 1-e + 2 \zeta, \cr
d & = & 1-e, \cr 
c & = & \frac12 \left (1 - \frac{e-\frac12}{\zeta} \right ). 
\end{eqnarray}
The degree function is exactly $\frac 12 - \zeta$ on the pode of size $c$ and $\frac12 + \zeta$ on the pode of size $1-c$.  

\begin{theorem}\label{thm:low-t} For sufficiently small $\zeta$, the entropy-maximizing graphon is 
unique and is well-approximated by the ansatz graphon $g_0$. Specifically, the 
entropy-maximing graphon has 
\begin{eqnarray}\label{what-is-g}
a & = & 1 - e - 2\zeta + O(\zeta^2), \cr 
b & = & 1-e + 2 \zeta + O(\zeta^2), \cr
d & = & 1-e + O(\zeta^2), \cr 
c & = & \frac12 \left (1 - \frac{e-\frac12}{\zeta} \right ) + O(\zeta).
\end{eqnarray}
Furthermore, the exact values of the parameters $a$, $b$, $c$, and $d$ 
are analytic functions of $(e,\tilde t)$ everywhere except at the singular point 
$(\frac12, 0)$. 
\end{theorem}

\begin{proof}[Proof of Theorem \ref{thm:low-t}]

Any graphon $g$ with degree function 
\begin{equation}
d(x) = \int_0^1 g(x,y) dy
\end{equation}
can be uniquely written as 
\begin{eqnarray}\label{general-g}
g(x,y) & = & d(x) + d(y) - e + \delta g(x,y) \cr 
& = & (d(x)-\frac12) + (d(y)-\frac12) + (1-e) + \delta g(x,y) \cr 
g(x,y)-\frac12 & = & (d(x)-\frac12) + (d(y)-\frac12) - (e - \frac12) + \delta g(x,y), 
\end{eqnarray}
where $\delta g(x,y)$ is a function with zero marginals: 
\begin{equation}
\int_0^1 g(x,y) \, dx = \int_0^1 g(x,y) \, dy = 0. 
\end{equation}
We will show that $\delta g(x,y)$ is pointwise $O(\zeta^2)$ and that $d(x) - \frac12$ only 
takes on two values, within $O(\zeta^2)$
of $\pm \zeta$. This implies that $g$ is bipodal and follows the estimates 
(\ref{what-is-g}). The analyticity of $(a,b,c,d)$ then follows from the implicit function
theorem. The proof follows 
several steps: 
\begin{enumerate}
\item Using known properties of the entropy function near Erd\H{o}s-R\'enyi to estimate the 
Lagrange multipliers $\alpha$ and $\beta$ that appear in equations (\ref{eq1}--\ref{eq3}).
\item Showing that $g$ is at most tripodal and that the degree function is everywhere 
$\frac12 + o(1)$. This implies that $g$ is pointwise close to $\frac12$.
\item Comparing the entropy of the general graphon $g(x,y)$ of equation (\ref{general-g})
to the entropy of the ansatz graphon $g_0$. This will show that $\iint \delta g(x,y)^2
dx \, dy = O(\zeta^6)$, that $g_0$ comes within $O(\zeta^6)$ of achieving the 
maximum possible entropy, and that the variance of $(d(x)-\frac12)^2$ is $O(\zeta^6)$. 
\item Since $(d(x)-\frac12)^2$ is almost constant (in an $L^2$ sense), and since 
$\int_0^1 (d(x)-\frac12)^2 \, dx = \zeta^2$ (exactly), our graphon must either be 
bipodal with degrees very close to $\frac12 \pm \zeta$, or must be tripodal with a 
very small third pode. We rule out the latter possibility. 
\item We examine the variational equations on the space of bipodal graphons in a 
neighborhood of $g_0$ and show that there is a unique solution that depends analytically
on $(e,\tilde t)$. 
\end{enumerate}

\medskip

\noindent {\bf Step 1:} The function $H(u)$
admits a convergent power series expansion around $u=\frac12$: 
\begin{equation}
H(u) = H\left (\frac12 \right ) + \frac12 \, H''\left ( \frac12 \right ) \left ( u-\frac12
\right )^2 + \frac{1}{24} \, H''''\left ( \frac12 \right ) \left (u-\frac12 \right )^4
+ \cdots.
\end{equation}
This gives rise to a convergent power series expansion for the entropy of a graphon $g$:
\begin{equation}
S(g) = \sum_{k=0}^\infty \frac{\mu_{2k}}{(2k)!} H^{(2k)}\left (\frac12\right ),
\end{equation}
where 
\begin{equation}
\mu_{2k} = \iint \left (g(x,y)-\frac12\right )^{2k} \,dx\,dy.
\end{equation}

When $\tilde t = 0$, the maximal entropy is exactly $H(e)$. As we vary $e$, the infinitesimal change in $t$ is $2e de$, so 
\begin{equation}
\alpha + 2e \beta = H'(e) \approx H''\left (\frac12 \right ) \left (e-\frac12\right ). 
\end{equation} 
The existence of the ansatz graphon $g_0$, with entropy $H(e) + H''(\frac12) \tilde t
+ O(\zeta^4)$, shows that $\beta$ is no less than 
$H''(\frac12) + O(\zeta) = -4+O(\zeta)$. However, if $\beta$ were greater than 
$H''(\frac12) + O(\zeta)$, then there would only be one solution to equation (\ref{eq3}),
namely $z=k(z)=e$. But that gives $\tilde t = 0$. When $\tilde t >0$, we must have 
\begin{equation}
\alpha = -H''(e) + O(\zeta), \qquad \beta = H''(e) + O(\zeta).
\end{equation}

\smallskip

\noindent{\bf Step 2:} With these values of $\alpha$ and $\beta$, the line $y=z$ is 
nearly tangent to $y=k(z)$ at $y=e$. This implies that all solutions to $z=k(z)$ are 
close to $e$, or equivalently close to $1/2$. 

The function $k(z)$ is the convolution of a fixed (scaled) logistic curve. The 
logistic function $1/(1+\exp(-(\alpha+\beta(z+e))))$ has a negative third derivative near $z=e$. Any small convolution
of this function must likewise have a negative third derivative, meaning that its 
second derivative is decreasing and only passes through zero once near $z=e$. 
By Rolle's theorem, this implies 
that a line can only intersect the graph $y=k(z)$ at most three times near $z=e$.
Thus our optimal graphon must be at most tripodal, with degrees close to 1/2. By
equation (\ref{eq1}), this means that the graphon is pointwise close to 1/2, meaning
that $g(x,y) - \frac12$ is $o(1)$ as $\zeta \to 0$. 

\smallskip

\noindent{\bf Step 3:} We now compute the leading terms in the expansions of $S(g_0)$ and $S(g)$. For any graphon, 
let 
\begin{equation}
\nu_k = \int_0^1 \left (d(x) - \frac12 \right)^k \, dx. 
\end{equation}
Note that the first two moments are determined by $e$ and $\tilde t$:
\begin{equation} \nu_1= e - \frac12 \hbox{ and } \nu_2 = \zeta^2 = \tilde t + \left ( e - \frac12 \right )^2. 
\end{equation} 
Higher even moments are bounded from below: 
\begin{equation} 
\nu_{2k} \ge \nu_2^k,
\end{equation}
with equality if and only if $\left( d(x)-\frac12 \right )^2$ is constant. In particular, $\nu_4-\nu_2^2$ is the variance of 
 $\left( d(x)-\frac12 \right )^2$.

 Let 
 \[ \eta = \left ( \iint \delta g(x,y)^2 \, dx], dy \right )^{1/2} \]
 be the $L^2$ norm of $\delta g$. For an arbitrary graphon, the first two non-trivial moments are:
 \begin{eqnarray}
 \mu_2 & = & 2 \nu_2 - \left ( e - \frac12\right )^2 + \eta^2 \cr 
 \mu_4 & = & 2 \nu_4 +  6\nu_2^2 - 12 \left (e-\frac12\right)^2 \nu_2 + 5\left (e-\frac12\right)^4 \cr 
 && + 24 \iint \delta g(x,y) \left [ (d(x)-\frac12)^2(d(y-\frac12)- (e-\frac12) (d(x)-\frac12)(d(y-\frac12)\right ]   dx \, dy \cr 
 && + 12 \iint \delta g(x,y)^2 \left [ (d(x)-\frac12)^2- (d(x)-\frac12)(d(y)-\frac12)\right ]   dx \, dy \cr
  && + 4 \iint \delta g(x,y)^3 \left [ 2(d(x)-\frac12)- (e-\frac12) \right ]   dx \, dy \cr
  && + \iint \delta g(x,y)^4  dx \, dy + 6 (e-\frac12)^2 \|\delta g\|_2^2. 
 \end{eqnarray}
 For the graphon $g_0$, this simplifies to 
  \begin{eqnarray}
 \mu_2 & = & 2 \nu_2 - \left ( e - \frac12\right )^2  \cr 
 \mu_4 & = & 2 \nu_2^2 -12 \nu_2 (e-\frac12)^2 + 5(e-\frac12)^4 \cr 
 & = & 2\zeta^4 + 6 (\zeta^2 - (e-\frac12)^2)^2 - (e-\frac12)^4. 
 \end{eqnarray}
Comparing these, we see that there is an $H''(1/2)\eta^2$ cost in having $\delta g$ nonzero and a $H''''(1/2) (\nu_4-\nu_2^2)/12$ 
cost in having $(d(x) - \frac12)^2$ not be constant. There are also costs in $\mu_4$ from terms proportional to $\delta g^2$ or $\delta g^4$. 

The lowest-order benefits from having $\delta g$ nonzero are terms proportional to 
\begin{eqnarray}
\iint \delta g(x,y) \left (d(x)-\frac12 \right )^2\left ( d(y)-\frac12 \right ) dx\, dy &&  \hbox{and} \cr 
\iint \delta g(x,y) (e - \frac12) \left (d(x)-\frac12 \right )\left ( d(y)-\frac12 \right ) dx\, dy && \end{eqnarray}
By Cauchy-Schwarz, the first term is $O(\zeta \eta \sqrt{\nu_4})$, while the second 
is $O(\zeta^3\eta)$. 
Since $\nu_4$ cannot be much greater than $\nu_2^2$, the maximum possible 
benefit is $O(\zeta^3\eta)$. 

Any benefits from the 
expansion of $\mu_6$ and higher are higher order in $\zeta$, $\eta$, or both, so the total benefit of having 
$\delta g$ nonzero is $O(\zeta^3\eta)$
With a cost proportional to 
$\eta^2$ and benefits that are $O(\zeta^3\eta)$, $\eta$ must itself be $O(\zeta^3)$ and the net benefit of having 
$\delta g$ nonzero is $O(\zeta^6)$. 

This means that the net cost associated with having $(d(x)-\frac12)^2$ differ from a constant must itself be $O(\zeta^6)$. 
There are indeed benefits at higher order, for instance terms proportional to $(e-\frac12) \nu_2\nu_3$ that appear in the
expansion of $\mu_6$, but they are all $O(\zeta^6)$, so the cost in $\mu_4$, proportional to the variance of 
$(d(x)-\frac12)^2$, must be $O(\zeta^6)$. 

\smallskip

\noindent{\bf Step 4:} We recall some facts about cubic polynomials. Suppose that 
\begin{equation}
f(x) = x^3 + a_2 x^2 + a_1 x + a_0 
\end{equation}
has three real roots. The sum of the roots is $-a_2$ and the average of the roots is $-a_2/3$, which is also the unique
point where $f''(x)=0$. Now consider the convolution of $f$ with a distribution of degree functions $d(y)$:
\begin{equation}
\tilde f(x) = \int_0^1 f(x+d(y)) dy.
\end{equation}
If $\tilde f$ has three roots, then the sum of the roots is $-3e-a_2$, where $e = \int_0^1 d(y) dy$ ,and the average of 
the roots is $-e - \frac{a_2}{3}$. Furthermore, a simple algebraic calculations shows that 
\begin{equation} 
\tilde f(-e - \frac{a_2}{3}) = f(-\frac{a_2}{3}) + \int_0^1 (d(y)-e)^3 \, dy.
\end{equation}

Now consider the function 
\[ \frac{1}{1 + \exp(-(\alpha + \beta z)} \]
Near the point of inflection $z=-\alpha/\beta$, this function is approximately cubic, with corrections of order 
$(z + \frac{\alpha}{\beta})^5$. The value of this function at the point of inflection is exactly 1/2. 
Convolving this function by the degree function $d(y)$ we obtain the function $k(z)$. The solutions to $k(z)$ have average
value $\bar d = -e - \frac{\alpha}{\beta}$. The value of $k(\bar d)$ is $\frac12 + \int_0^1 (d(x)-e)^3 \, dx$, plus
$O((\bar d-e)^5)$ corrections due to $k(z)$ not being exactly cubic. 

We have already shown that two of the three roots of $k(z)-z$ are $\frac12 \pm \zeta + O(\zeta^2)$. Since 
$\int_0^1 d(x)^3 \, dx$ is then $O(\zeta^3)$, this implies that the average of the three roots is $\frac12 + O(\zeta^3)$,
which implies that the third root must be $\frac12 + O(\zeta^2)$. Let $c_3$ be the size of this pode. 

We now compute the cost 
\begin{equation} 
\nu_4 - \nu_2^2 = \int_0^1 ((d(x)-\frac12)^2-\zeta^2)^2 \, dx \approx c_3 \zeta^4. 
\end{equation} 
The derivative of the entropy with respect to $c_3$ has a positive term of order $\zeta^4$. Since $d(x)-\frac12$ is pointwise
$O(\zeta)$, all of the contributions to $\mu_6$ and higher have order $\zeta^6$ and higher, and cannot overcome this 
quartic cost. Nor can the cross-terms with $\delta g$, which we have shown to be $O(\zeta^6)$. Since the derivative of the 
entropy with respect to $c_3$ is positive, the entropy is maximized when $c_3=0$. That is, the optimal graphon is 
bipodal, not tripodal, with the degree functions on the two podes being $\frac12 \pm \zeta + O(\zeta^2)$. 

\smallskip

\noindent{\bf Step 5:} Bipodal graphons are described by four parameters $(a,b,c,d)$. The edge density, 2star density, and 
entropy are all analytic functions of $(a,b,c,d)$. Introducing Lagrange multipliers, we obtain four 
analytic equations in six unknowns: 
\begin{equation}
\nabla S  =  \alpha \nabla e + \beta \nabla t 
\end{equation}
This gives a 2-dimensional analytic variety of solutions. To see that $(a,b,c,d)$ are analytic functions of $e$ and $t$,
we need only check that the tangent space does not degenerate. That is, we must check that on this family of solutions, 
the edge and triangle densities can be varied independently to first order. 

However, that is easy. This property obviously holds for the ansatz graphon (\ref{eq:ansatz}), except at the singularity 
$\zeta=0$ (where the derivative of $\zeta$ with respect to $\tilde t$ diverges). 
$\partial d/\partial e=-1$, while $\partial d/\partial \tilde t=0$. However, 
the partial derivatives of $a$, $b$ and $c$ with respect to $\tilde t$ are nonzero, showing
that $\partial_e(a,b,c,d)$ and $\partial_{\tilde t}(a,b,c,d)$ are linearly independent 
vectors. In particular, 
\begin{equation} 
\frac{\partial a}{\partial e} \frac{\partial d}{\partial \tilde t} - 
\frac{\partial d}{\partial e} \frac{\partial a}{\partial \tilde t} = \frac{1}{\zeta}.
\end{equation}

The difference between the true graphon and the ansatz is small, in particular being
$O(\zeta^2)$ for $a$ and $d$, with the derivatives of these terms with respect to 
$\zeta$ being $O(\zeta)$. Since 
\begin{equation}
\frac{\partial \zeta}{\partial e} = \frac{e-\frac12}{\zeta}=O(1) \hbox{ and }
\frac{\partial \zeta}{\partial \tilde t} = \frac{1}{2\zeta},
\end{equation}
these terms can only change $\partial_e(a)$ and $\partial_e(d)$ by $O(\zeta)$ and 
$\partial_{\tilde t}(a)$ and $\partial_{\tilde t}d$ by $O(1)$, resulting in an $O(1)$
change in $\partial_ea \partial_{\tilde t}d - \partial_e d \partial_{\tilde t}a$,
which remains nonzero for small $\zeta$. 

\end{proof}

\end{document}
